\subsection{Ordered Sets of Submodules}

In this subsection, A and B are k-algebras, M is a left A-module, and V is a left B-module. We denote by \(\mathrm{D}(\mathrm{M})\) (resp. \(\mathrm{D}(\mathrm{V})\) ) the set of submodules of M (resp. of V), ordered by inclusion. We assume given an isomorphism of ordered sets \(\varphi : \mathrm{D}(\mathrm{V}) \to \mathrm{D}(\mathrm{M})\) .

By Morita's theorem (VIII, p.~103, Theorem 2), we obtain such an isomorphism in the following situation: P is an invertible \((\mathrm{A},\mathrm{B})_k\) -bimodule, M is the A-module \(\mathrm{P} \otimes_{\mathrm{B}} \mathrm{V}\) , and for every submodule W of V, \(\varphi(\mathrm{W})\) is the canonical image of \(\mathrm{P} \otimes_{\mathrm{B}} \mathrm{W}\) in M.

Some properties of the module M, or of its submodules, can be expressed in terms of the ordered set D(M): they are listed in Tables I and II.

The module M is the direct sum of a family \((\mathrm{M}_{i})_{i\in\mathrm{I}}\) of submodules if and only if we have \(M=\sum_{i\in I}M_{i}\) and \(M_{i}\cap\sum_{j\neq i}M_{j}=0\) for every \(i\in I\) . This remark and an examination of Table I give the following result.

PROPOSITION 7. --- a) We have \(\varphi(0)=0\) and \(\varphi(V)=M\) .

\begin{enumerate}
\def\labelenumi{\alph{enumi})}
\setcounter{enumi}{1}
\tightlist
\item
  Let \((\mathrm{V}_i)_{i\in \mathrm{I}}\) be a family of submodules of V. We have
\end{enumerate}

\[
\varphi \left(\sum_ {i \in \mathrm{I}} \mathrm{V} _ {i}\right) = \sum_ {i \in \mathrm{I}} \varphi (\mathrm{V} _ {i}), \quad \varphi \left(\bigcap_ {i \in \mathrm{I}} \mathrm{V} _ {i}\right) = \bigcap_ {i \in \mathrm{I}} \varphi (\mathrm{V} _ {i}).
\]

Submodules of M

Ordered set D(M)

Zero submodule

Smallest element of D(M)

Submodule M

Greatest element of D(M)

\(\bigcap_{i\in I} M_i\)

Greatest lower bound \(\inf_{i\in I} M_i\)

\(\sum_{i\in I} M_i\)

Least upper bound \(\sup_{i\in I} M_i\)

Supplementary submodules

\(\inf(M', M'') = 0, \sup(M', M'') = M\)

Simple submodule of M

Minimal element of D(M) --- \{0\}

Maximal submodule of M

Maximal element of D(M) --- \{M\}

Socle \(\mathscr{S}(M)\) of M

Least upper bound in D(M) of the set of minimal elements of D(M) --- \{0\}

\emph{Radical \(\Re(M)\) of M}(VIII, p.~151)

Greatest lower bound in D(M) of the set of maximal elements of D(M) --- \{M\}

Properties of the module M

Properties of D(M)

M is Noetherian.

The ordered set D(M) is Noetherian (Set Theory, III, §6, No.~5, p.~190).

M is Artinian.

The set D(M), ordered by ⊃, is Noetherian.

M is indecomposable.

We have M ≠ 0, and there are no two nonzero elements M\textquotesingle{} and M\textquotesingle\textquotesingle{} of D(M) satisfying inf(M\textquotesingle, M\textquotesingle\textquotesingle) = 0, sup(M\textquotesingle, M\textquotesingle\textquotesingle) = M.

M is finitely generated.

For every family (Mi) i∈I in D(M) with upper bound M, there exists a finite subset J of I such that M = supj∈I Mj.

M is simple.

Card(D(M)) = 2.

M is semisimple.

M is the least upper bound, in D(M), of the set of minimal elements of D(M) = \{0\}.

TABLE II.

\begin{enumerate}
\def\labelenumi{\alph{enumi})}
\setcounter{enumi}{2}
\tightlist
\item
  The B-module V is the direct sum of the family \((\mathrm{V}_{i})_{i\in\mathrm{I}}\) of submodules if and only if M is the direct sum of the family \((\varphi(\mathrm{V}_{i}))_{i\in\mathrm{I}}\) .
\end{enumerate}

Let \(V'\) and \(V''\) be submodules of \(V\) such that \(V'\) is contained in \(V''\) ; set \(M' = \varphi(V')\) and \(M'' = \varphi(V'')\) , so that \(M''\) contains \(M'\) . Denote by \([V', V'']\) the interval in \(D(V)\) consisting of the submodules \(W\) of \(V\) such that we have \(V' \subset W \subset V''\) , and define the interval \([M', M'']\) in \(D(M)\) likewise. The mapping \(W \mapsto W/V'\) is an isomorphism of ordered sets from \([V', V'']\) to \(D(V''/V')\) ; we define an isomorphism of ordered sets from \([M', M'']\) to \(D(M''/M')\) likewise. Since \(\varphi\) sends the interval \([V', V'']\) to \([M', M'']\) , it defines an isomorphism \(\varphi\) of ordered sets from \(D(V''/V')\) to \(D(M''/M')\) . We deduce the following proposition from this and Tables I and II.

PROPOSITION 8. --- a) Let V' and V'\,' be submodules of V such that V'\,' contains V'. The B-module V'\,'/V' is simple if and only if the A-module \(\varphi(\mathrm{V}^{\prime\prime})/\varphi(\mathrm{V}^{\prime})\) is simple.

\begin{enumerate}
\def\labelenumi{\alph{enumi})}
\setcounter{enumi}{1}
\item
  If \(V'\) is a simple submodule, maximal submodule, or direct factor of V, then \(\varphi(V')\) is, respectively, a simple submodule, maximal submodule, or direct factor of M.
\item
  The isomorphism \(\varphi\) transforms the socle \(\mathcal{S}(\mathrm{V})\) of \(\mathrm{V}\) into the socle \(\mathcal{S}(\mathrm{M})\) of \(\mathrm{M}^*\) and the radical (VIII, p.~151, Definition 1) \(\Re (\mathrm{V})\) of \(\mathrm{V}\) into the radical \(\Re (\mathrm{M})\) of \(\mathrm{M}_{*}\)
\item
  Let \((\mathrm{V}_{i})_{0\leqslant i\leqslant n}\) be a finite sequence of submodules of V. It is a Jordan-Hölder series of V if and only if \((\varphi(\mathrm{V}_{i})_{0\leqslant i\leqslant n})\) is a Jordan-Hölder series of M.
\end{enumerate}

Lemma 2. --- Let H and H' be submodules of V such that \(H \cap H' = 0\) . The B-modules H and H' are isomorphic if and only if the A-modules \(\varphi(\mathrm{H})\) and \(\varphi(\mathrm{H}')\) are isomorphic.

We identify \(H+H'\) with the product \(H \times H'\) . The graph of an isomorphism from H to \(H'\) is a submodule \(H''\) of V satisfying

\[
\mathrm{H} \cap \mathrm{H} ^ {\prime \prime} = \mathrm{H} ^ {\prime} \cap \mathrm{H} ^ {\prime \prime} = 0, \quad \mathrm{H} + \mathrm{H} ^ {\prime} = \mathrm{H} + \mathrm{H} ^ {\prime \prime} = \mathrm{H} ^ {\prime} + \mathrm{H} ^ {\prime \prime};\tag{17}
\]

conversely, every submodule that has these properties is the graph of an isomorphism from H to \(H'\) . By Proposition 7, the relation \(H \cap H' = 0\) is equivalent to \(\varphi(\mathrm{H}) \cap \varphi(\mathrm{H}') = 0\) , and relations (17) are equivalent to the relations

\[
\varphi (\mathrm{H}) \cap \varphi (\mathrm{H} ^ {\prime \prime}) = \varphi (\mathrm{H} ^ {\prime}) \cap \varphi (\mathrm{H} ^ {\prime \prime}) = 0,
\]

\[
\varphi (\mathrm{H}) + \varphi (\mathrm{H} ^ {\prime}) = \varphi (\mathrm{H}) + \varphi (\mathrm{H} ^ {\prime \prime}) = \varphi (\mathrm{H} ^ {\prime}) + \varphi (\mathrm{H} ^ {\prime \prime});
\]

the lemma follows.

PROPOSITION 9. --- Let S be a simple submodule of V and T the simple submodule \(\varphi(\mathrm{S})\) of M. If \(V_{S}\) denotes the isotypical component of V of type S and \(M_{T}\) the isotypical component of M of type T, then we have \(\varphi(V_{\mathrm{S}})=M_{\mathrm{T}}\) .

Every simple submodule \(S'\) of V distinct from S satisfies \(S' \cap S = 0\) . It is therefore isomorphic to S if and only if \(\varphi(S')\) is isomorphic to T (Lemma 2). Now, \(V_{S}\) is the sum of the simple submodules of V isomorphic to S, and \(M_{T}\) is the sum of the simple submodules of M isomorphic to T. Proposition 9 therefore follows immediately from Propositions 7 and 8.

PROPOSITION 10. --- a) The B-module V is Artinian (resp. Noetherian, indecomposable, simple, finitely generated) if and only if M is.

\begin{enumerate}
\def\labelenumi{\alph{enumi})}
\setcounter{enumi}{1}
\item
  The B-module V has finite length if and only if the A-module M has finite length, and we then have \(\mathrm{long}_{\mathrm{B}}(\mathrm{V}) = \mathrm{long}_{\mathrm{A}}(\mathrm{M})\) .
\item
  The B-module V is semisimple (resp. isotypical) if and only if the A-module M is semisimple (resp. isotypical). If this is the case, then we have \(\text{long}_{\text{B}}(\text{V}) = \text{long}_{\text{A}}(\text{M})\) .
\end{enumerate}

Assertion a) follows from the second property listed in Table II.

Assertion b) follows from Proposition 8, d).

The module V is semisimple if and only if it is equal to its socle \(\mathcal{S}(V)\) ; it is isotypical if and only if there exists a simple submodule S of V such that \(V = V_{S}\) . Assertion c) therefore follows from Propositions 7, c), 8, c), and 9 (VIII, p.~106 and 109).

